% Abhakseite „Das kann ich“ – Zone nur im Gesamt (\mitzone)
%% TODO Ich-kann-Sätze: Platzhalter wie in den Aufgabendateien
\begin{abhakseite}
\ifdefined\mitzone
\abhakgruppe{Kennst du schon}
\abhak{1}{Multiplizieren mit Vorzeichen und Punkt vor Strich}
\abhak{2}{Termwert berechnen, auch mit negativer Zahl, als Probe einer Umformung}
\abhak{3}{Gleichartige Glieder mit einer Variablen zusammenfassen, Vorzahl eins beachten, Potenz getrennt halten}
\abhak{4}{Zahl mal Klammer und Minusklammer (Distributivgesetz, alle Vorzeichen drehen)}
\abhak{5}{Quadratzahlen bis 15² erkennen und Quadratwurzeln daraus}
\abhak{6}{Scheitelpunktform lesen}
\abhak{7}{Ausklammern eines gemeinsamen Zahlfaktors oder einer Variablen}
\abhak{8}{Zahl mal Klammer und Minusklammer (Distributivgesetz, alle Vorzeichen drehen) – Fehler finden}
\abhak{9}{Zahl mal Klammer und Minusklammer (Distributivgesetz, alle Vorzeichen drehen) – selbst rechnen}
\fi
\abhakgruppe{Einheit 1 · Summe mal Summe}
\abhak{10}{Summe mal Summe}
\abhak{11}{Termwert mit zwei Variablen berechnen (auch negativ)}
\abhak{12}{Summe mal Summe – Fehler finden, Begründen}
\abhakgruppe{Einheit 2 · Binomische Formeln}
\abhak{13}{Gleiche Klammer zweimal?}
\abhak{14}{Binomische Formeln}
\abhak{15}{Kopfrechnen mit der Formel (Vorrat)}
\abhak{16}{Binomische Formeln – Fehler finden, Begründen}
\abhakgruppe{Einheit 3 · Faktorisieren}
\abhak{17}{Faktorisieren}
\abhak{18}{Faktorisieren – Fehler finden, Begründen}
\end{abhakseite}
